\section{Uniform exponential integrability}

Skoda's local theorem gives an integrable exponential for each quasi-plurisubharmonic potential. Normalized compactness makes the exponent and integral bound uniform: otherwise a convergent weighted sum of counterexamples would violate the local theorem.

Write $K_\omega$ for the bound in \cref{thm:normalized-compactness}, so that normalized potentials $v\in\PSH(X,\omega)$ with $\sup_Xv=0$ satisfy
\begin{equation}\label{eq:normalized-lone}
 \int_X|v|\,\mathrm{d}V\leq K_\omega.
\end{equation}

Local Skoda integrability implies that every quasi-plurisubharmonic function $v$ on compact $X$ admits some exponent $a(v)>0$ for which $\mathrm{e}^{-a(v)v}$ is integrable. Indeed, apply the local theorem to a plurisubharmonic correction near each point, and then take the minimum exponent over a finite cover. We use the local integrability theorem recalled in \cite{SGPT}.

\begin{lemma}\label{lem:uniform-exponential}
There are constants $\alpha>0$ and $C<+\infty$, depending on $(X,\omega,\mathrm{d}V)$, such that
\[
 \int_X\mathrm{e}^{-\alpha v}\,\mathrm{d}V\leq C
 \quad\hbox{for all }v\in\PSH(X,\omega),\quad \sup_Xv=0.
\]
\end{lemma}
\begin{proof}
Suppose no such pair exists. Set $a_j=2^{-j}$ and choose normalized potentials $v_j\leq0$ with
\[
 \int_X\exp\!\left(-\frac{a_j}{j}v_j\right)\mathrm{d}V\geq2^j.
\]
The partial sums $V_N=\sum_{j\leq N}a_jv_j$ decrease, are $\omega$-plurisubharmonic, and have uniformly bounded $L^1$ norms by \eqref{eq:normalized-lone}. Their limit $V$ is therefore $\omega$-plurisubharmonic. Since all summands are nonpositive, $V\leq a_jv_j$ for every $j$.

\Needspace{6\baselineskip}
Fix $p>0$. For every $j\geq1/p$,
\[
 \int_X\mathrm{e}^{-pV}\,\mathrm{d}V
 \geq\int_X\mathrm{e}^{-pa_jv_j}\,\mathrm{d}V
 \geq\int_X\mathrm{e}^{-(a_j/j)v_j}\,\mathrm{d}V
 \geq2^j.
\]
Thus $\mathrm{e}^{-pV}$ is nonintegrable for every $p>0$, contradicting the qualitative Skoda input applied to $V$.
\end{proof}
